% 本文件由 LaTeX 排版 Agent 根据有效 translation.json 自动生成。
% author=Jerome; work_id=3008; title=首位隐修士保禄传

\chapter{首位隐修士\Person{保禄}传}

% structure: p0001 source=h1 target=omitted reason=page-title-already-rendered-by-parent

\WorkTitle{首位隐修士保禄传}写于374或375年，即\Person{热罗尼莫}在\Place{叙利亚}旷野期间（见第6章），并献给\Place{康科迪亚}的\Person{保禄}，如\Person{热罗尼莫}的\WorkTitle{书信}第十篇第三章所述。\par

1. 关于哪位隐修士最早树立隐修生活的典范，这一问题广泛且频繁地被讨论。有些人追溯过远，在圣人们\Person{厄里亚}和\Person{若翰}身上找到了开端——前者似乎不仅是一位隐修士，后者在出生前就已开始预言。另一些人，他们的观点普遍被接受，认为\Person{安当}是这种生活方式的创始者，这种看法部分正确。我说部分，因为事实并非他领先于众人，而是众人都从他那里获得了必要的激励。但直到今日，\Person{安当}的两位门徒\Person{阿玛塔斯}和\Person{玛加略}（前者将他的老师安葬）断言，一位\Place{特拜}的\Person{保禄}（虽非第一个以保禄为名的人）是运动的先驱，这一看法也得到我的赞同。有些人随意编造故事，说他是一个住在地下的洞穴里的人，头发垂到脚上，还编造了许多难以置信的传说，这些细节不值一述。而且，毫无羞耻地撒谎的人的观点也不值得驳斥。既然\Place{希腊}和\Place{罗马}的作家都已仔细记载了\Person{安当}的事迹，我决定简略记述\Person{保禄}的早年与晚年，更多是因为此事被忽略了，而非出于对自己能力的自信。他的中年生活如何，以及他经历了\Person{撒殚}的何种陷阱，据认为尚无人发现。\par

2. 在\Person{德西乌斯}和\Person{瓦勒里安}迫害期间，当\Place{罗马}的\Person{科尔乃略}和\Place{迦太基}的\Person{西彼廉}在蒙福的殉道中倾流鲜血时，\Place{埃及}和\Place{特拜}地区的许多教会都被这场风暴的愤怒所摧毁。那时，基督徒常祈求能为\Person{基督}之名死于刀剑。但狡猾的敌人的愿望是杀害灵魂，而非身体；他通过仔细寻找缓慢而致命的酷刑来达到目的。用亲身遭受其手的\Person{西彼廉}本人的话来说：\Quote{那些想死的人不被允许被杀}。我们举两个例子，两者都特别值得注意，并能使敌人的残酷更加为人所知。\par

3. 一位在信德上坚贞不移的殉道者，在刑架和烙板之间屹立如得胜者，敌人命令他全身涂满蜂蜜，双手反绑在背后，躺在烈日之下，这样，已经克服了煎锅之热的人，可能被苍蝇的叮刺所征服。另一人正值青春年华，被敌人命令带到某个宜人的游乐花园里，那里有洁白的百合和羞红的玫瑰，旁边是潺潺的小溪，头顶上微风的轻柔低语在树叶间嬉戏，他被放在一张铺满羽毛的深床之上，用甜美的花环捆住以防他逃脱。当所有人都退去后，一位极美的娼妓走近，开始以淫荡的拥抱将手臂缠在他的脖子上，并且——那可耻地提及！——抚摸他的身体，以便一旦肉体的情欲被激起，她就可以实现她的淫荡目的。\Person{基督}的士兵不知该做什么，往何处去。不为酷刑所屈的他，正在被快乐征服。最后，他受天上的灵感，咬断了舌头的尖端，在她亲吻时吐在她脸上。这样，情欲的感觉就被随之而来的剧痛所制服。\par

4. 当这样的暴行在特拜地区下部分肆虐时，\Person{保禄}和他刚结婚的妹妹失去了双亲，他当时约十六岁。他继承了一笔丰厚的遗产，精通\Place{希腊}和\Place{埃及}的学问，生性温柔，并深爱天主。在迫害的雷声中，他退到一处相当远的、更隐秘的房屋中。但\Quote{被诅咒的贪金之欲}将人的心引向何等罪恶！他的妹夫起了念头，要出卖这个他本应庇护的青年。妻子的眼泪——它们常常能打动人心——、血缘的纽带，以及上方天主无所不见的眼睛，都不能使这叛徒脱离他的邪恶。\Quote{他来了，他紧逼，他残忍行事，表面却只像在强调亲情。}\par

5. 年轻人有智慧理解这一点，并顺服意志于必然，逃往山区旷野，等待迫害结束。他以轻松的行程和多次的停留，开始向沙漠深处行进。最后他发现一座石山，山脚下有一块石头封住的一个不大的洞穴。他移开石头（人多么急切地想知道隐藏之物），急切地搜寻，看到里面有一个大厅，露天但被一棵古棕榈的宽大枝叶遮蔽。棕榈并没有遮掩一泉清澈的泉水，泉水刚涌出就被同一地面的一个小开口吞没。此外，在充满洞穴的山中，还有许多可居住的地方，里面可以看到锈迹斑斑的冲压钱币的砧子和锤子。\Place{埃及}作家记载，这地方在\Person{安当}与\Person{克娄巴特拉}结合时是一个秘密铸币厂。\par

6. 于是，他将这居所视为天主的恩赐，爱上了它，并在祈祷和独居中度过余生。棕榈为他提供了食物和衣服。为了使无人认为这不可能，我以\Person{耶稣}及其圣天使为证，我在\Place{叙利亚}与萨拉森人之间那片沙漠中见过，并且至今仍见，一些隐修士，其中一人被关了三十年，靠吃大麦面包和泥水维生；另一人在一个旧蓄水池（在\Place{叙利亚}方言中称为\OriginalTerm{古巴}{Gubba}）中，每天靠五个干无花果维持生命。我所叙述的如此奇特，对于不相信\BibleQuote{为信的人，一切都是可能的}这话的人，会显得难以置信。\par

7. 但让我们回到我岔开的话题。真福保禄在地上已度了一百一十三年天堂般的生活，而安东尼在九十岁时正住在另一处僻静之地（他本人常这样说），这时他心中起了个念头：没有比自己更成全的隐修士留在旷野了。然而，在夜晚的寂静中，他得到启示：在旷野深处还有一位比他好得多的人，他应该去拜访他。于是，天一亮，这位可敬的长者便拄着拐杖支撑并引导自己衰弱的四肢出发了；但他不知道该朝哪个方向走。灼热的正午到了，头顶烈日炙烤，但他并未因此放弃已开始的旅程。他说：\Quote{我信赖我的天主：总有一天祂会将祂应许给我的同仆指示给我。}他不再多说。忽然，他看见一个混合形状的生物，半马半人，诗人们称之为半人马（\OriginalTerm{半人马}{Hippocentaur}）。一见到它，他就在额头上划出救恩的标记来武装自己，然后喊道：\Quote{喂！这一带哪里有天主的仆人生存？}那怪物先是发出一些古怪的、嘶哑的、与其说是说话不如说是从毛茸茸的嘴唇里挤出的不成词句的声音，终于找到了友善的沟通方式，伸出右手指出了所寻求的道路。然后它以迅疾的飞驰穿过了辽阔的平原，从惊讶的同伴眼前消失了。但这究竟是魔鬼变成这形状来恐吓他，还是旷野因以盛产奇异动物而也生出这种生物，我们无法断定。\par

8. 安东尼大为惊奇，思考着所见之事，继续上路。不久，在一个四面环山的小岩石谷中，他看见一个小矮人，钩鼻、额生角、脚像山羊蹄。安东尼见状，像一位好战士一样拿起信德的盾牌和望德的神盔；那生灵却开始向他献上棕榈树果实，供他旅途之用，作为和平的保证。安东尼看出这一点，停下脚步问它是谁。它回答他说：\Quote{我是有死的受造物，是那些被外邦人以各种错误形态迷惑、称之为法翁（\OriginalTerm{法翁}{Fauns}）、萨提尔（\OriginalTerm{萨提尔}{Satyrs}）和因库比（\OriginalTerm{因库比}{Incubi}）的旷野居民之一。我是代表我的族类被派来的。我恳求你为我们求取你我的主——我们得知祂曾来拯救世界，\BibleQuote{祂的声音传遍了普世}（\BibleRef{罗 10:18}）——的恩宠。}当它说出这样的话时，老旅人的双颊流下了泪水，那是他深情的标志，他在满心的喜乐中洒下泪水。他为基督的光荣和撒殚的覆灭而欢喜，同时惊叹于自己能听懂萨提尔的语言。他用拐杖敲击地面，说：\Quote{祸哉，亚历山大里亚！你敬拜怪物，却轻慢天主！祸哉，你这娼妓之城，全世界的魔鬼都涌入了你！如今你还有什么话说？野兽讲论基督，而你却敬拜怪物，轻慢天主。}他话还没说完，那野物便如长了翅膀般逃走了。但愿无人对这件事心生疑虑；它的真实性得了君士坦丁皇帝在位时发生之事的支持，那件事整个世人都曾目睹。因为一个那样的活人被带到亚历山大里亚，作为奇景展示给民众。之后，为防止尸体在夏季腐烂，它被盐腌保存，带到安提约基雅让皇帝观看。\par

9. 继续讲述我计划中的故事。安东尼穿越他所进入的地区，只看见野兽的踪迹和旷野的广阔荒芜。他不知道该做什么，该往何处去。又过去了一天。他只剩下了一样东西：坚信自己不会被基督抛弃。第二个夜晚的黑暗在祈祷中度过了。天刚蒙蒙亮，他看见不远处有一只母狼因干渴而喘着气，爬向山脚。他用目光追随它；那野兽消失在一个洞穴中后，他走近并向内探望。他的好奇心毫无所获：黑暗阻碍了视线。但正如圣经所说：\BibleQuote{完美的爱德驱散恐惧}（\BibleRef{若一 4:18}）。他蹒跚着脚步，屏住呼吸，小心摸索着走了进去；他一点一点前进，反复倾听声音。最后，在恐怖的午夜黑暗中，远处出现了亮光。他急切地加快步伐，脚踢到一块石头，引起了回声；于是真福保禄关上了开着的门，用闩子拴住。安东尼便倒在入口处的地上，直到第六时或更晚，他恳求准入，说：\Quote{我是谁、从哪里来、为什么来，祢知道。我知道自己不配看见祢，但除非见到祢，我绝不离开。祢欢迎野兽，为何不欢迎人？我求了，我就找到；我敲门，就给我开门。但若我不成功，我就死在祢门槛上。我死后，祢必定会埋葬我。}\par

他这样不停地呼喊，却屹然不动。\newline{}英雄对他这样简短地回答：\par

\Quote{这样的祈祷不是威胁；眼泪中没有诡计。祢来到这里求死，我不欢迎祢，祢奇怪吗？}保禄微笑着让他进入，门开了，他们彼此拥抱，互相叫出名字，一同感谢天主。\par

10. 在神圣的亲吻之后，\Person{保禄}坐下，这样开始对\Person{安东尼}说：\Quote{看，这就是你费尽辛劳寻找的人，他的肢体因年迈而衰败，他的白发蓬乱。你看到面前这个人，他不久将成为尘土。但是爱德忍受一切。因此我求你告诉我，人类的情况如何？古城中是否有新的家园兴起？什么政府治理世界？是否还有一些人留下来被魔鬼的幻象带走？}他们这样交谈时，惊奇地注意到一只乌鸦落在树枝上，然后轻轻飞下，来到他们面前放下一整块面包。他们惊讶，乌鸦飞走后，\Quote{看，}\Person{保禄}说，\Quote{主真是仁慈，真是慈爱，给我们送来了食物。过去六十年，我总是收到半块面包；但在你来时，基督加增了他战士的口粮。}\par

11. 相应地，感谢主之后，他们一起坐在晶莹的泉边。此时，关于谁该擘饼发生了争论，几乎一整天直到黄昏都在讨论。\Person{保禄}以款待为由支持自己的观点，\Person{安东尼}则提出年龄。最后，他们商定每人抓住靠近自己一边的饼，拉向自己，并保留自己手中的部分。然后他们用手和膝盖喝了一点泉水，并向天主献上赞颂之祭，整夜醒着祈祷。天亮时，真福\Person{保禄}这样对\Person{安东尼}说：\Quote{兄弟，我早知道你住在那地方；很久以前天主就应许把你作为我的同仆给我；但现在我安眠的时刻近了；我一直渴望被解脱而与基督同在；我的赛程已完，正义的冠冕为我存留。因此，主派你来将我可怜的尸体安葬在土里，是将尘土归于尘土。}\par

12. \Person{安东尼}听到这话，流着泪，呻吟着，开始祈求\Person{保禄}不要离开他，而要让他作为旅伴同行。\Person{保禄}回答说：\Quote{你不该寻求自己的益处，而该寻求别人的益处。对你而言，卸下肉身的负担，跟随羔羊是相宜的；但对于其余的弟兄们，由你的榜样来训练是相宜的。所以，劳驾你回去取\Person{亚他那修}主教给你的那件外套，来包裹我可怜的身体。}真福\Person{保禄}提出这个请求，并非因为他很在乎他的尸体腐朽时是否有衣服遮盖（既然他很久以来一直穿着用棕榈叶缝制的衣服，他何必在乎呢？），而是为了减轻朋友对他去世的悲伤。\Person{安东尼}惊讶地发现\Person{保禄}听说了\Person{亚他那修}和他的外套；他好像在其中看见基督自己，心里朝拜天主，不敢多说一句话；然后默默流泪，再次亲吻他的眼睛和双手，就出发返回后来被撒拉逊人占领的隐修院。他的脚步跟不上他的意愿。然而，他虽然因禁食而精疲力竭，又被年迈削弱，但他的勇气战胜了他的年龄。\par

13. 最后，他疲惫不堪、气喘吁吁地完成了行程，到达了他的小屋。这里两个门徒前来迎接他，他们在他年老时开始服侍他。他们说：\Quote{父啊，你在哪里待了这么久？}他回答说：\Quote{我罪人何其不幸！我不配称为隐修士。我看见了\Person{厄里亚}，我看见了沙漠中的\Person{若翰}，我真的看见了乐园中的\Person{保禄}。}然后他闭上嘴，捶着胸，从房间取出外套。当门徒们请他更详细地解释时，他说：\Quote{有缄默的时候，也有说话的时候。}\BibleRef{训 3:7}\par

14. 然后他出去了，没有吃一点东西，就按原路返回，只渴望着他，渴望见到他，眼里心里只有他。因为他担心，后来事件证明他的预感是正确的，怕在他不在时，他的朋友会把自己的灵魂交付给基督。又一个黎明来了，还有三小时的路程，他看见\Person{保禄}身穿雪白的衣袍，在天使的群中上升，还有先知和宗徒的队列。他立即俯伏在地，把粗沙撒在头上，痛哭流涕地喊道：\Quote{\Person{保禄}，你为什么抛弃我？为什么不告别一声就走？你为何这么晚才让我认识你，却又这么快就离去？}\par

15. 真福\Person{安东尼}后来常叙述说，他在剩下的路程中走得如此之快，像鸟一样飞；这不无原因：因为进入洞穴时，他看到尸体呈跪姿，头抬起，双手举起。他做的第一件事是以为他还活着，就在他旁边祈祷。但是当他没有听到祈祷者通常发出的叹息时，他就亲吻哭泣，然后他明白了：即使圣人的尸体也以虔敬的姿势向天主祈祷，万物在主内都是活的。\par

16. 随后，他包好遗体并抬出去，一路按基督徒传统咏唱圣歌和圣咏，\Person{安当}便开始哀叹没有挖掘土地的器具。于是，在思绪的汹涌大海中，他思量了许多计划，说道：若我返回隐修院，那是四日的路程；若我留在这里，也无济于事。那么，就让我像理应如此一样，在你这位战士身旁死去，噢，\Person{基督}，我将很快咽下最后一口气。正当他反复思量这些事时，看，两头鬃毛飞扬的狮子从旷野深处疾驰而来。起初他见此情景惊恐万分，但再次将思绪转向\Person{天主}后，他毫无惊慌地等待着，仿佛他看到的是鸽子。它们径直来到那位有福老人的尸体前，停住，向他摇尾乞怜，并在他的脚边躺下，大声吼叫，仿佛是要表明它们以唯一可能的方式哀悼。然后它们开始在近旁刨地，竞相挖掘沙子，直到挖出一个恰好容下一人的地方。随即，仿佛要求为它们的工作得到奖赏，它们竖起耳朵，垂下头，来到\Person{安当}跟前，开始舔他的手和脚。他明白它们是在向他求祝福，立刻向\Person{基督}发出赞美，称连哑巴动物也感受到祂的天主性，他说：\Quote{主，没有你的命令，叶子不会从树上落下，麻雀不会掉到地上，请赐予它们你所知道的最好的恩惠。}然后他挥手示意它们离开。它们走后，他弯下年迈的肩膀承受圣人的遗体，放入坟墓，用挖出的土覆盖，并在上面堆起惯常的坟头。另一天破晓，为了让这位充满感情的后嗣不缺少那位没有遗嘱的死者的一些东西，他为自己取下了圣人用棕榈叶像编织品一样织成的长衣。然后返回隐修院，他向他所有的门徒详细讲述了这一切，并在复活节和五旬节的庆日上总是穿着\Person{保禄}的长衣。\par

17. 请允许我在这篇小论的结尾问那些不知道自身财产限度的人，那些用大理石装饰房屋的人，那些把房屋连成片、把田地连成片的人：这位赤身的老人曾缺少过什么吗？你们的酒杯是宝石做的，他用掌心来解渴；你们的长衣是金线织的，他连你们最卑微的奴隶的衣物也没有。然而，他虽贫穷，天堂却为他敞开；你们虽富有金银，却将被收入地狱。他赤身露体，却保存了基督的衣裳；你们身着绸缎，却失去了基督的衣服。\Person{保禄}躺在无用的尘土中，但将复活进入荣耀；你们有昂贵的坟墓被修建起来，但你们和你们的财富都注定被焚烧。我恳求你们，至少，请为你们所爱的财富有所顾虑。为何连你们死者的殓衣也用金子制成？为何在哀悼和泪水中你们的夸耀也不停止？难道富人的尸体只有丝绸才能腐烂吗？\par

18. 我恳求你，读者，无论你是谁，请记得罪人\Person{热罗尼莫}。他，如果\Person{天主}给他选择，他宁愿取\Person{保禄}的长衣连同他的功德，也不取君王的紫袍连同他们的惩罚。\par

\SourceRecord{Translated by W.H. Fremantle, G. Lewis and W.G. Martley.；Nicene and Post-Nicene Fathers, Second Series；Vol. 6.；Edited by Philip Schaff and Henry Wace.；Buffalo, NY: Christian Literature Publishing Co.,；1893.；<http://www.newadvent.org/fathers/3008.htm>.}
